\documentclass[12pt]{article}
\usepackage[amsmath]{e-jc}
\usepackage{booktabs}
\usepackage[T1]{fontenc}
\usepackage{lmodern}
\usepackage{microtype}
\usepackage{listings}
\lstset{basicstyle=\small\ttfamily,breaklines=true,columns=fullflexible}
\newcommand{\B}{\beta_*}
\newcommand{\cost}{\mathcal Q}
\newcommand{\E}{\mathbb E}
\newcommand{\ind}{\mathbf 1}
\title{Sparse-set extension and degree-sensitive\\
perturbations for halves of triangle-free graphs}
\author{komeiji shiki}
\hypersetup{pdfauthor={komeiji shiki}}
\MSC{05C35}
\Copyright{The author.}
% Initial-submission metadata only. The unmodified journal style controls
% typography. Publication dates, volume, article number and DOI are assigned
% by the editors after acceptance. Replace this block with their supplied
% \dateline and \specs values when preparing an accepted final version.
\makeatletter
\def\the@dateline{Manuscript prepared for submission, 5 September 2026\\[0.5ex]\copy@right}
\def\ps@plain{\let\@oddhead\@empty\let\@evenhead\@empty
  \def\@oddfoot{\hfil\thepage\hfil}\let\@evenfoot\@oddfoot}
\makeatother
\pagestyle{plain}
\begin{document}
\maketitle
\begin{abstract}
Erd\H{o}s's sparse half conjecture asks whether every triangle-free graph
on $n$ vertices contains $\lfloor n/2\rfloor$ vertices spanning at most
$n^2/50$ edges. We give a computer-assisted proof of the general bound
$(2543/100000)n^2=0.02543n^2$. The proof combines degree-sensitive
neighbourhood perturbations, extension from a sparse vertex set, and
neighbourhoods anchored across a maximum cut. A reciprocal-degree averaging
criterion reduces one construction to an inequality in three real
variables. Five exact rational subdivision certificates cover its entire
parameter domain. The remaining ingredients are known maximum-cut bounds
and a four-cycle density certificate, which we check with a separate
implementation enumerating every triangle-free graph on six vertices.
\end{abstract}

\section{Statement and relation to previous work}
For a finite simple graph $G$ on $n\ge1$ vertices, put
\[
 b(G)=\min_{|S|=\lfloor n/2\rfloor}e(G[S]),\qquad
 \rho(G)=\frac{2e(G)}{n^2}.
\]
The sparse half conjecture asks whether $b(G)\le n^2/50$ whenever $G$ is
triangle-free; see the problem entry \cite{erdos}. Razborov \cite{razborov}
proved the general bound $27n^2/1024$. Sarid \cite{sarid} subsequently gave
a preprint claiming $131n^2/5000$, with a rational four-cycle certificate
and a Lean development conditional on a quoted max-cut theorem. We use
the numerical four-cycle certificate from that work, but neither import
its verifier nor rely on its Lean development.

\begin{theorem}\label{thm:main}
Every finite triangle-free graph $G$ on $n\ge1$ vertices satisfies
\[
 b(G)\le qn^2,\qquad q=\frac{2543}{100000}=0.02543.
\]
\end{theorem}

The numerical ingredients are a finite coefficient check and five
interval checks. The first checks an existing four-cycle certificate;
the others check a scalar inequality covering a continuum of parameters.
None is
a search restricted to graphs of bounded order. In particular, the earlier
finite graph census in the accompanying research repository is not an
assumption of Theorem~\ref{thm:main}.

The degree-dependent endpoint radii themselves already occur in
\cite{sarid}. Lemma~\ref{lem:tangent} uses their dependence on the degree
in a reciprocal-degree averaging criterion.
Lemma~\ref{lem:extension} allows the starting set to contain edges.
Section~\ref{sec:opposite} uses neighbourhoods from the other side of
the cut. We combine these with the density-specific max-cut bound.
Neither the final constant nor the parameter $19/20$ is claimed optimal.

\section{Fractional halves and the external ingredients}
Throughout the proof, $G$ is finite, simple and triangle-free.
Define
\[
 \B(G)=\frac1{n^2}\min\left\{
 \sum_{\{u,v\}\in E(G)}f_uf_v:
 0\le f_v\le1,\quad\sum_v f_v=n/2\right\}.
\]
Then $b(G)/n^2\le\B(G)$. Indeed, among minimizers choose one with as few
fractional coordinates as possible. Moving two such coordinates as
$(f_u+t,f_v-t)$ preserves their sum. The objective is affine or concave
quadratic in $t$, so one endpoint of the feasible interval is no worse.
Thus at most one coordinate is fractional. For even $n$ all are integral;
for odd $n$, discard the sole possible coordinate of value $1/2$.
We will use that $\B(G(k))=\B(G)$ for every balanced blow-up $G(k)$.
Indeed, replace a profile on each twin class by its average. The mass
constraint scales by $k$ and every edge product scales by $k^2$;
conversely a profile on $G$ lifts constantly to its twin classes.

Write $D_2(G)$ for the least number of edges lying within the two parts of
a bipartition. We use the following consequences of Theorem~2(a,b) in the
PDF version of \cite{bcl}:
\begin{equation}\label{eq:cuts}
 \frac{D_2(G)}{n^2}\le\frac2{47},\qquad
 \rho(G)\ge\rho_0:=\frac{3197}{10000}
 \ \Longrightarrow\ \frac{D_2(G)}{n^2}\le\frac1{25}.
\end{equation}
That reference states its results for sufficiently large graphs, and the
density condition there is $e(G)\ge0.3197\binom n2$. To obtain
\eqref{eq:cuts}, replace each vertex of $G$ by $k$ nonadjacent twins.
The resulting graph $G(k)$ is triangle-free and
$D_2(G(k))=k^2D_2(G)$. For the nontrivial inequality in this identity,
let $x_i$ be the fraction of class $i$ on one side of a cut in $G(k)$.
Its internal edge count divided by $k^2$ is
\[
 \sum_{\{i,j\}\in E(G)}\bigl(x_ix_j+(1-x_i)(1-x_j)\bigr).
\]
This is affine in each individual coordinate, so its minimum over the
cube occurs at an integral point, corresponding to a cut of $G$.
Also $\rho(G(k))=\rho(G)$, and $\rho(G)\ge0.3197$ implies
$e(G(k))\ge0.3197\binom{nk}{2}$. Take $k$ sufficiently large and divide
the cited bounds by $k^2$.

Our second input is the certificate from \cite{sarid}:
\begin{equation}\label{eq:c4}
 t(C_4,G)\ge\rho(G)^3-\gamma\rho(G),\qquad
 \gamma=\frac3{64}+\frac1{50\,000\,000\,000}.
\end{equation}
Here $t(C_4,G)=\operatorname{hom}(C_4,G)/n^4$, with noninjective maps
included. Appendix~\ref{app:c4} describes the independent coefficient
check and the precise sampling argument turning it into this inequality.
The max-cut theorem \eqref{eq:cuts} is cited, not reproved by our programs.

We will also use the neighbourhood estimate
\begin{equation}\label{eq:anchor}
 \B(G)\le\frac\rho8-\frac{t(C_4,G)}{4\rho}\quad(\rho>0),
\end{equation}
which appears with a different four-cycle normalization in \cite{razborov}.
For completeness, if a neighbourhood has size at least $n/2$, an
independent fractional half has cost zero. The right side of
\eqref{eq:anchor} is nonnegative: the edge-counting identity below and
$e(N(u),N(v))\le e(G)$ imply $t(C_4,G)\le\rho^2/2$.
Otherwise every neighbourhood has size less than $n/2$. Fix an edge $uv$,
set $A=N(u)$, $C=N(v)$ and
$B=V(G)\setminus(A\cup C)$, and put
\[
 p=\frac{1/2-|A|/n}{1-|A|/n-|C|/n}.
\]
The sets $A,C$ are disjoint and independent, and $p\in[0,1]$.
The fractional halves $\ind_A+p\ind_B$ and
$\ind_C+(1-p)\ind_B$ have costs $Q_u,Q_v$ satisfying
\[
 (1-p)Q_u+pQ_v
 =p(1-p)\left(\frac\rho2-\frac{e(A,C)}{n^2}\right)
 \le\frac14\left(\frac\rho2-\frac{e(A,C)}{n^2}\right).
\]
Average over all edges, using
$\sum_{uv\in E(G)}e(N(u),N(v))=\operatorname{hom}(C_4,G)/2$,
to obtain \eqref{eq:anchor} whenever the zero-cost case does not apply.

\section{Keeping the degree in the perturbation}
Let $H$ be a triangle-free graph with positive vertex weights $w_x$
summing to one. Define
\[
 \mu=\sum_{xy\in E(H)}w_xw_y,\quad
 d_x=\sum_{y\in N(x)}w_y,\quad
 e_x=\sum_{y\in N(x)}w_yd_y,\quad
 S_x=d_x(e_x-\mu d_x).
\]
Assume $\mu>0$. Let $\sigma$ denote the total weight of nonisolated
vertices. Cauchy--Schwarz in $N(x)$ gives
\[
 \frac{d_x^2}{e_x}\le\sum_{y\in N(x)}\frac{w_y}{d_y}.
\]
All the denominators are positive. Summing with weights $w_x$ and
reversing the order of summation proves
\begin{equation}\label{eq:reciprocal}
 \sum_{x:d_x>0}w_x\frac{d_x^2}{e_x}\le\sigma,
 \qquad \sum_xw_xd_x=2\mu.
\end{equation}

For $0<a\le1/2$ and $0<d<1$, let
\begin{equation}\label{eq:radii}
 r(a,d)=
 \min\left\{\frac{1-a}{1-d},\frac a d\right\}
 \min\left\{\frac a{1-d},\frac{1-a}d\right\}.
\end{equation}

\begin{lemma}\label{lem:tangent}
Fix $0<a\le1/2$ and $K>0$. Suppose that, for every $d\in(0,1)$ such that
$19d>18\mu$,
\begin{equation}\label{eq:scalar-condition}
 K\le r(a,d)\mu d^2\frac{d+18\mu}{19d-18\mu}.
\end{equation}
Then some nonisolated vertex $x$ satisfies $r(a,d_x)S_x\ge K$.
Consequently, for every linear functional $\ell$, the constant profile
$a$ on $H$ can be replaced by a profile $\pi\in[0,1]^{V(H)}$ of the same
weighted mean such that
\[
 \cost(\pi)+\ell(\pi)\le\cost(a)+\ell(a)-K,
 \qquad \cost(\pi)=\sum_{xy\in E(H)}w_xw_y\pi_x\pi_y.
\]
The same assertion holds with constant profile $1-a$.
\end{lemma}

\begin{proof}
Suppose $r(a,d_x)S_x<K$ for all nonisolated $x$. Then
\[
 \frac{d_x^2}{e_x}>
 \frac{d_x^3}{\mu d_x^2+K/r(a,d_x)}.
\]
For $19d_x\le18\mu$, the right side is positive and hence at least
$19d_x/(20\mu)-9/10$. In the other case, this same lower bound follows
by rearranging \eqref{eq:scalar-condition}, with positive denominators.
Summing the resulting strict inequalities yields
\[
 \sum_{x:d_x>0}w_x\frac{d_x^2}{e_x}
 >\frac{19}{10}-\frac9{10}\sigma\ge\sigma,
\]
contradicting \eqref{eq:reciprocal}.

Choose the vertex just obtained and abbreviate $d=d_x$,
$\varphi=\ind_{N(x)}-d$. Its weighted mean is zero. As $N(x)$ is
independent, expansion of the quadratic form gives
\[
 \cost(\varphi)=-d(e_x-\mu d)=-S_x.
\]
The feasible interval for $a+t\varphi$ is $-t_-\le t\le t_+$, where
\[
 t_+=\min\{(1-a)/(1-d),a/d\},\quad
 t_-=\min\{a/(1-d),(1-a)/d\}.
\]
Give its two endpoints probabilities $t_-/(t_++t_-)$ and
$t_+/(t_++t_-)$. Then $\E t=0$ and $\E t^2=t_+t_-=r(a,d)$.
All linear terms cancel, including $\ell$, and the expected change in
cost is $-r(a,d)S_x\le-K$. At least one endpoint has that improvement.
Replacing $a$ by $1-a$ exchanges $t_+$ and $t_-$, leaving their product
unchanged.
\end{proof}

The coefficient $19/20$ is a chosen tangent parameter. More generally one
can compare the reciprocal expression with
$\lambda d/\mu+1-2\lambda$. Our choice allows the verification to use
only rational functions; it is not claimed to be optimal.

\section{Extending a sparse set}\label{sec:extension}
The following argument adapts the greedy extension and neighbourhood
sampling in Section 4.5 of \cite{razborov}. Its starting set need not be
independent. For $1/3\le z\le1/2$ put
\[
 J(z)=\begin{cases}9/128-z/8,&z\le3/8,\\
 z(1/2-z)/2,&z\ge3/8.\end{cases}
\]
For $1/3\le A\le3/8$ define
\begin{align*}
 K_A(z)&=\begin{cases}
 (12A-32z^2-12z+9)/192,&z\le3/8,\\
 (A-8z^2+3z)/16,&z\ge3/8,
 \end{cases}\\
 \Phi_A(z)&=\max\{J(z),K_A(z)\}.
\end{align*}
Both $J$ and $K_A$ are continuous and decreasing on this interval.

\begin{lemma}\label{lem:extension}
Let $G$ be triangle-free. Suppose its normalized independence number is at most $A$,
where $A\in[1/3,3/8]$ is rational. If $Z\subseteq V(G)$ has size $zn$
with $1/3\le z\le1/2$ and spans $\varepsilon n^2$ edges, then
\[
 \B(G)\le\varepsilon+\Phi_A(z).
\]
\end{lemma}
\begin{proof}
Use a balanced blow-up to assume that $n/8$ and $An$ are integers;
all normalized parameters, including the independence number, are
unchanged. For the latter claim, an independent set in the blow-up meets
only classes whose original vertices are pairwise nonadjacent; taking
all vertices of those classes cannot decrease its size.
Starting at $B=Z$, repeatedly add a vertex outside $B$ with
at most $n/8$ neighbours in $B$, stopping on reaching $n/2$ vertices or
when none is available. Write $b=|B|/n$ and $t=1/8$. Throughout,
\[
 \frac{e(B)}{n^2}\le\varepsilon+t(b-z).
\]
If $b=1/2$, this is at most $\varepsilon+J(z)$. Suppose instead the
algorithm stops early. Let $C$ be the vertices outside $B$ with more
than $bn/2$ neighbours in $B$. Any two vertices of $C$ have a common
neighbour in $B$, so $C$ is independent and $|C|\le An$.
Choose $D$ disjoint from $B\cup C$ with normalized size
$D_0=1-A-b$. Every $w\in D$ has normalized degree into $B$ in
$(t,b/2]$. Such a choice exists because $|C|/n\le A$.
Put $r=1/2-b$; in particular, $D_0-r=1/2-A\ge1/8$.

If some $v\in B$ has at least $rn$ neighbours in $D$, choose $rn$ of
them as $E$. This is an independent set, and $B\cup E$ has normalized
cost at most
\[
 \varepsilon+t(b-z)+\frac b2(1/2-b)\le\varepsilon+J(z).
\]
The last quadratic has its maximum at $b=3/8$ when $z\le3/8$,
and at $b=z$ otherwise.

It remains to consider the case that every $v\in B$ has fewer than
$rn$ neighbours in $D$. Choose $v$ uniformly in $B$, let
$E=N(v)\cap D$, $F=D\setminus E$, and $s_v=|E|/n\le r$.
Take the fractional half equal to $1$ on $B\cup E$ and to
$(r-s_v)/(D_0-s_v)$ on $F$. The latter ratio lies in $[0,1]$,
since $D_0>r$. Bounding $e(B,F)/n^2$ by $b|F|/(2n)$,
those between $E,F$ trivially, and the edges inside $F$ by Mantel's
bound gives its cost at most
\[
 \frac{e(B)+e(B,E)}{n^2}
 +\frac14(r-s_v)(1/2+b+3s_v).
\]
Let $s=\E s_v=e(B,D)/(bn^2)$, so $s\le\min\{r,D_0/2\}$.
For $u\in D$, write $h_u=|N(u)\cap B|/n\in[t,b/2]$.
The inequality $h_u^2\le(t+b/2)h_u-tb/2$ and double counting give
\[
 \E\frac{e(B,E)}{n^2}
 =\frac1{bn}\sum_{u\in D}h_u^2
 \le(t+b/2)s-\frac{tD_0}{2}.
\]
The other term is concave in $s_v$, so Jensen's inequality bounds the
expected cost by $\varepsilon$ plus
\begin{equation}\label{eq:extension-quadratic}
 t(b-z)+\frac{r(1-r)}4+(r/2+t)s-\frac34s^2-\frac{tD_0}{2}.
\end{equation}
For $b\le3/8$, the unconstrained maximum in $s$ is at
$s=(r+1/4)/3$, which is no greater than either $r$ or $D_0/2$
because $b\le3/8\le3/2-3A$. For $b\ge3/8$, one has
$r\le D_0/2$ since $b\ge A$, and the maximum on
$0\le s\le\min\{r,D_0/2\}$ is at $s=r$.
After substitution, the two expressions are
\[
 \frac{12A-32b^2+12b-24z+9}{192},\qquad
 \frac{A-8b^2+5b-2z}{16},
\]
respectively. They coincide at $b=3/8$. Their derivatives in $b$ are
$-(16b-3)/48$ and $-(16b-5)/16$, negative on their relevant intervals
with $b\ge z\ge1/3$. Thus the maximum occurs at $b=z$ and equals
$K_A(z)$. One of the sampled halves is no more expensive than its
expectation. This proves the lemma.
\end{proof}

\begin{lemma}\label{lem:independence}
If a triangle-free graph $G$ has an independent set of normalized size at least $9/25$,
then $\B(G)\le81/3200<2543/100000$.
\end{lemma}
\begin{proof}
Let $\alpha$ be its normalized independence number. For
$9/25\le\alpha\le3/8$, apply Lemma~\ref{lem:extension} to a maximum
independent set with $A=z=\alpha$ and $\varepsilon=0$.
Here $K_\alpha(\alpha)=3/64-\alpha^2/6$, and
\[
 J(\alpha)-K_\alpha(\alpha)=(\alpha-3/8)^2/6\ge0.
\]
Hence $\B(G)\le J(\alpha)\le J(9/25)=81/3200$.
For $\alpha\ge3/8$, Theorem 3.6 of \cite{razborov} gives
$\B(G)\le\alpha(1/2-\alpha)/2\le3/128$ when
$\alpha\le1/2$; the case $\alpha\ge1/2$ is immediate.
\end{proof}

\section{Neighbourhoods anchored across a maximum cut}\label{sec:opposite}
Choose a cut $X,Y$ realizing $D_2(G)$, with $|X|\le|Y|$, and retain
$c=n/(2|Y|)$, $a=1-c$, $\mu=e(Y)/|Y|^2$. Suppose the independence
number is at most $An$, where $A=9/25$. Define
\[
 h=\frac{|X|}{|Y|}=2c-1,\quad T=2Ac,\quad k=\min\{T,h\},
 \quad Z=\frac{e(X,Y)}{|Y|^2}.
\]
We use this construction only when $a\le1/4$, so $h>0$.
For $u\in Y$, put $d_u=|N(u)\cap Y|/|Y|$ and
$\psi_u=|N(u)\cap X|/|Y|$. Optimality of the cut under a single
vertex flip implies $\psi_u\ge d_u$. Since every neighbourhood in a
triangle-free graph is independent, the total degree is at most $An$.
Consequently
\[
 0\le d_u\le T/2,\quad d_u\le\psi_u\le k,
 \quad\E_Yd_u=2\mu,\quad\E_Y\psi_u=Z.
\]
From $(\psi_u-d_u)(k-\psi_u)\ge0$ and Cauchy--Schwarz,
\begin{equation}\label{eq:product-moment}
 \E_Y(d_u\psi_u)\ge
 \E_Y\frac{k d_u^2}{k-\psi_u+d_u}
 \ge\frac{4k\mu^2}{k-Z+2\mu}.
\end{equation}
Zero terms with a zero denominator are interpreted by continuity, or
omitted before applying Cauchy--Schwarz. The last denominator is positive
when $\mu>0$, since $Z\le k$.

For $x\in X$, let $I_x=N(x)\cap Y$, an independent set in $G[Y]$,
and put
\[
 D_x=\frac{|I_x|}{|Y|},\quad
 E_x=\frac1{|Y|}\sum_{u\in I_x}d_u,
 \quad S_x=D_x(E_x-\mu D_x).
\]
Double counting gives
\[
 \E_XD_x=Z/h,\qquad
 \E_XE_x=\E_Y(d_u\psi_u)/h.
\]
Moreover, $0\le D_x\le T$ and
$0\le E_x\le\min\{\mu,TD_x/2\}$. The first bound on $E_x$ uses
the independence of $I_x$: each edge incident to $I_x$ is counted once.

\begin{lemma}\label{lem:opposite}
With the maximum-cut notation and independence cap above, assume
$0<a\le1/4$, $\mu>0$, $\mu\le T^2/2$, and set
\[
 v=\mu(1-T)-2\mu^2/T,\qquad
 W=\frac{4Tk\mu^2}{h(k-Z+2\mu)}+\frac{vZ}{h}-T\mu+2\mu^2.
\]
For a constant profile $a$ or $1-a$ on $Y$, with any added linear
functional, a perturbation within $Y$ saves
at least $L(a)\max\{W,0\}/(4c^2)$ in ambient normalization, where
$L(a)=\min\{4a^2,a(1-a)\}$.
If $v\ge0$, this conclusion remains true after replacing $Z$ in $W$
by any nonnegative lower bound for $Z$.
\end{lemma}
\begin{proof}
For $0\le D\le T$ and $0\le E\le\min\{\mu,TD/2\}$,
\begin{equation}\label{eq:curvature-plane}
 D(E-\mu D)\ge TE+vD-T\mu+2\mu^2.
\end{equation}
Indeed, the difference has coefficient $D-T\le0$ in $E$, so it
suffices to check its upper endpoint. Put $D_*=2\mu/T\le T$.
For $D\le D_*$ the difference equals
$(T-D)(T/2-\mu)(D_*-D)$; for $D\ge D_*$ it equals
$\mu(T-D)(D-D_*)$. Both are nonnegative, since $T\le1$ and
$\mu\le T^2/2$ imply $\mu\le T/2$.
Average \eqref{eq:curvature-plane} and apply
\eqref{eq:product-moment} to obtain $\E_XS_x\ge W$.
Thus some anchor has $S_x\ge W$.

The perturbation $\ind_{I_x}-D_x$ has mean zero and quadratic cost
$-S_x$, exactly as in Lemma~\ref{lem:tangent}. The product of its two
feasible radii is $r(a,D_x)\ge L(a)$ when $0<D_x<1$.
If $W>0$, the selected anchor has $D_x>0$; also $D_x\le T\le18/25<1$.
To see the radius inequality,
use the three ranges $D_x\le a$, $a\le D_x\le1-a$, and
$D_x\ge1-a$: their products are respectively
$a(1-a)/(1-D_x)^2$, $a^2/(D_x(1-D_x))$, and
$a(1-a)/D_x^2$. If $W>0$, choose the better endpoint; if $W\le0$,
keep the original profile. Finally divide the saving by $4c^2$.
For $v\ge0$, both $Z$-dependent terms in $W$ are increasing in $Z$
on $0\le Z\le k$, proving the last assertion.
For a substituted lower bound $Z_0\le Z$, one has
$\E_X S_x\ge W(Z)\ge W(Z_0)$, so the same anchor-selection argument
applies to $\max\{W(Z_0),0\}$.
\end{proof}

\section{The cut construction and the continuous check}
Take a cut $V(G)=X\sqcup Y$ realizing $D_2(G)$, with $|X|\le|Y|$ and
$I=(e(X)+e(Y))/n^2\le I_0$. Write
\[
 c=1-a=\frac{n}{2|Y|},\quad 0\le a\le\frac12,
 \qquad \mu=\frac{e(Y)}{|Y|^2}\le4I_0c^2.
\]
The profiles $\ind_X+a\ind_Y$ and $c\ind_Y$ are fractional halves.
Their costs are respectively
\begin{equation}\label{eq:baselines}
 B_1=\frac{a\rho}2+cI-\frac{1+a}{4c}\mu,
 \qquad B_2=\frac\mu4.
\end{equation}
Inside $Y$, a change to either profile contributes a quadratic cost in
$G[Y]$ and an arbitrary linear functional coming from edges to $X$.
The conversion from the normalization on $Y$ to that on $G$ multiplies
costs by $|Y|^2/n^2=1/(4c^2)$.

By Lemma~\ref{lem:independence} it remains to treat graphs with
independence number less than $(9/25)n$. If $a\le1/4$, then
$z=|X|/n=1-1/(2c)\ge1/3$, and Lemma~\ref{lem:extension} gives
the additional bound
\begin{equation}\label{eq:extension-candidate}
 \B(G)\le E(a,\mu):=I_0-\frac{\mu}{4c^2}
      +\Phi_{9/25}\left(1-\frac1{2c}\right).
\end{equation}
For $a>1/4$, regard this candidate as unavailable.
When a lower bound $\rho\ge R_-$ is also specified, use
$Z_0=4c^2(R_-/2-I_0)\le Z$ in Lemma~\ref{lem:opposite}, provided
its hypotheses and $Z_0\ge0$, $v\ge0$ hold. With $B(a,\mu)$ defined
below, this supplies the auxiliary candidate
\[
 O(a,\mu)=B(a,\mu)-\frac{L(a)}{4c^2}\max\{W(Z_0),0\}.
\]
An inapplicable auxiliary candidate is simply omitted.
The lower bound for $Z$ follows from the exact identity
$Z=4c^2(\rho/2-I)$.

The following scalar check combines the candidates.
\begin{lemma}\label{lem:interval}
Let $q=2543/100000$ and take $(R_-,R,I_0)$ from the density table
in Appendix~\ref{app:interval}. Define
\[
 P(a)=I_0+a(R/2-I_0),\qquad
 B(a,\mu)=\min\left\{P(a)-\frac{1+a}{4(1-a)}\mu,\frac\mu4\right\}.
\]
Whenever $0<a\le1/2$, $0<\mu\le4I_0(1-a)^2$, $0<d<1$, and
$19d>18\mu$, either an available auxiliary candidate is at most
$q$, or one has
\begin{equation}\label{eq:interval}
 B(a,\mu)-
 \frac{r(a,d)\mu d^2(d+18\mu)}{4(1-a)^2(19d-18\mu)}\le q.
\end{equation}
\end{lemma}
\begin{proof}
Appendix~\ref{app:interval} gives rational upper bounds for the left side
on rectangular boxes and a finite subdivision covering the entire domain.
The accompanying checker verifies all leaves and the coverage of the
subdivision. All acceptance comparisons use rational arithmetic.
\end{proof}

\begin{proposition}\label{prop:low}
If $G$ is triangle-free and $\rho(G)\le7/20$, then $\B(G)\le q$.
\end{proposition}
\begin{proof}
Apply Lemma~\ref{lem:independence} first. Otherwise the independence
cap required in \eqref{eq:extension-candidate} holds. Use the row of
the density table containing $\rho$; the cuts exist by \eqref{eq:cuts}.
For $a=0$, the two parts are equal and one has cost at most $I_0/2<q$.
For $\mu=0$, the second baseline is zero. Thus assume both are positive.
Replacing $\rho,I$ in $B_1$ by $R,I_0$ can only increase it.
If $B(a,\mu)\le q$, a baseline suffices; if an available auxiliary
candidate is at most $q$, its construction suffices. Otherwise put
$K=4c^2(B(a,\mu)-q)>0$. Lemma~\ref{lem:interval} supplies
\eqref{eq:scalar-condition}. Apply Lemma~\ref{lem:tangent} within $Y$
to the better baseline. The resulting fractional half has cost at most
$B(a,\mu)-K/(4c^2)=q$.
\end{proof}

\section{The remaining densities}
For $\rho\ge7/20$, either a zero-cost half exists or
\eqref{eq:c4} and \eqref{eq:anchor} give
\[
 \B(G)\le\frac\rho8-\frac{\rho^2}{4}+\frac\gamma4
 \le\frac{159}{6400}+\frac1{200\,000\,000\,000}<q.
\]
The second inequality follows because the quadratic is decreasing for
$\rho\ge7/20$. Combining this with Proposition~\ref{prop:low} and the
rounding argument proves Theorem~\ref{thm:main}.

\appendix
\section{An exact certificate for the continuous inequality}\label{app:interval}
The five rows below cover all densities from zero through $7/20$.
In the first and last rows the opposite-side candidate is not used.
The final column gives total tree nodes, not sampled parameter points.
\begin{center}
\begin{tabular}{cccr}
\toprule
$R_-$ & $R$ & $I_0$ & Nodes\\
\midrule
$0$ & $3/10$ & $2/47$ & 5,049\\
$3/10$ & $31/100$ & $2/47$ & 445\\
$31/100$ & $63/200$ & $2/47$ & 1,273\\
$63/200$ & $3197/10000$ & $2/47$ & 1,331\\
$3197/10000$ & $7/20$ & $1/25$ & 467\\
\bottomrule
\end{tabular}
\end{center}
The starting box is
\[
 [0,1/2]\times[0,4I_0]\times[0,1]
\]
in the coordinates $(a,\mu,d)$. Let a subbox have endpoints
$[a_-,a_+]$, $[m_-,m_+]$, $[d_-,d_+]$, and put
$c_-=1-a_+$, $c_+=1-a_-$. A box is discarded if
$m_->4I_0c_+^2$, or if $19d_+\le18m_-$. In boxes with $a_+\le1/4$,
the extension candidate has the upper bound
\[
 U_E=I_0-\frac{m_-}{4c_+^2}
       +\Phi_{9/25}\left(1-\frac1{2c_-}\right).
\]
This uses the fact that $\Phi_{9/25}$ is decreasing in its argument.
Such a box is certified when $U_E\le q$. Otherwise the baseline is at most
\[
 U_B=\min\left\{\frac{m_+}{4},\
 I_0+a_+(R/2-I_0)-\frac{(1+a_-)m_-}{4c_+}\right\}.
\]
If $U_B\le q$, the whole box is certified. For the opposite-side
candidate, set $A=9/25$ and
\[
 T_\pm=2Ac_\pm,\quad h_\pm=2c_\pm-1,\quad
 k_\pm=\min(T_\pm,h_\pm),\quad
 Z_\pm=4c_\pm^2(R_-/2-I_0).
\]
Use this candidate only if the row enables it and
\[
 a_+\le1/4,\quad R_-\ge2I_0,\quad m_->0,\quad
 m_+\le T_-^2/2,\quad Z_+\le k_-,
\]
and
\[
 v_-=m_-(1-T_+)-2m_+^2/T_-\ge0.
\]
These conditions justify replacing $Z$ by $Z_0$ and ensure that all
compared denominators are positive. A lower bound for $W(Z_0)$ is
\[
 W_-=
 \frac{4T_-k_-m_-^2}{h_+(k_+-Z_-+2m_+)}
 +\frac{v_-Z_-}{h_+}-m_+T_++2m_-^2.
\]
Thus a box is also certified when
\[
 U_B-\frac{L(a_-)}{4c_+^2}\max(W_-,0)\le q.
\]
This comparison does not depend on the degree coordinate $d$.
Here $L$ is increasing on $[0,1/4]$: both $4a^2$ and $a(1-a)$
are increasing on that interval.

In the remaining boxes define
\[
 \mathcal R_-=
 \min\left\{\frac{c_-}{1-d_-},\frac{a_-}{d_+}\right\}
 \min\left\{\frac{a_-}{1-d_-},\frac{c_-}{d_+}\right\}.
\]
Every factor is a lower bound for its corresponding endpoint radius.
The denominator $19d_+-18m_-$ is positive in the remaining boxes.
An upper bound for the left side of \eqref{eq:interval} is therefore
\begin{equation}\label{eq:box}
 U_B-
 \frac{\mathcal R_-m_-d_-^2(d_-+18m_-)}{4c_+^2(19d_+-18m_-)}.
\end{equation}
Endpoint zeros give a valid zero lower bound for the saving. The domains
of interest have $0<d<1$; no surviving nondegenerate box has $d_-=1$
or $d_+=0$.

The certificate is a preorder binary tree. Symbols $0,1,2$ bisect the
corresponding coordinate at its rational midpoint; $L$ denotes a leaf.
Both children are checked and their union equals their parent, including
the shared boundary. The checker rejects any missing nodes, trailing
nodes, invalid symbols, or leaves failing all the stated tests.
Thus the certificate covers a continuum, including all boundaries,
rather than a discrete sampling of points.

\begin{center}
\begin{tabular}{lr}
\toprule
Certificate statistic & Combined count\\
\midrule
Total tree nodes & 8,565\\
Leaves certified by the baseline & 327\\
Leaves certified by sparse-set extension & 841\\
Leaves certified by opposite-side anchors & 736\\
Leaves with nonpositive tangent & 27\\
Leaves certified by the degree tangent & 2,036\\
Leaves outside the mass constraint & 318\\
Maximum tree depth & 28\\
\bottomrule
\end{tabular}
\end{center}
The verification command below prints the certificate file names.
In the order of the density table, their hashes are
\begin{lstlisting}
8a5aa7beda22b8133ce48b0556cac51ee77558cbb4dc370f9a5dba9f17d934ec
f4b4819fc870a0b5e715b7d0c77544091aa6131cc1bca6af4b2680778b70f855
2d755294dbbbdfac6c92a555001870caa25c393ba651fe89d40f795779fe28a6
b8dbe23d14b0c3ceb260cae9f4da0772c1804004c0db31d27e6eb2ac2670a5eb
f8ad90fdb9d0934d91f6f45831c71de7bddc97ae8a309fdd4a177a4fb91c1ee6
\end{lstlisting}
Replay the stored certificates with Python 3.11:
\begin{lstlisting}
python -S src/verify_paper.py
\end{lstlisting}
These programs need only the Python standard library. A floating-point
heuristic is used only to choose which coordinate to split when generating
the tree; it has no role in accepting a leaf. The checker does not use that
heuristic. This is a reproducible computer-assisted check, not a formal
proof-assistant verification of the Python interpreter or the reduction.
A second implementation, \texttt{src/audit\_intervals.py}, imports no
other project modules. It uses generic rational interval operations and
recursive tree traversal and checks the same 8,565 nodes independently
of the first implementation. It can be run separately:
\begin{lstlisting}
python -S src/audit_intervals.py
\end{lstlisting}
Its report records exact minimum margins for each leaf type in every row.

The optional script \path{src/audit_algebra.py} uses SymPy 1.14.0 to expand
12 polynomial identities in the proof exactly; these checks are not
numerical samples and are not required by the standard-library replay.

\section{Independent verification of the four-cycle certificate}\label{app:c4}
The numerical certificate belongs to Amir Sarid \cite{sarid}, at commit
\begin{lstlisting}
ba85f319ebb33ec2194ef4c2929a12916329e1e7
\end{lstlisting}
It is retained as \texttt{certificates/sarid-clebsch-tangent.json}, with
SHA-256
\begin{lstlisting}
606b26d835d5b7fe70d077230f8091bbd9677283049a1d0b46a3e3ad30a200ce
\end{lstlisting}
Our separate implementation enumerates all $2^{15}=32768$ labelled
six-vertex simple graphs. Exactly 5,789 are triangle-free. Removing their
full relabelling orbits gives 38 representatives, without relying on an
external catalogue. The program rederives the coordinate orderings for
the certificate's flag vectors and checks all 34 positive rational
rank-one terms against all 38 hosts. Every coefficient slack is
nonnegative; they are recorded in \texttt{results/four-cycle-independent.json}.

Here is why this finite check implies \eqref{eq:c4} for every finite $G$.
Sample six vertices of $G$ independently with replacement and let $H$ be
the graph of adjacencies on their six positions. It is a triangle-free
simple graph, even when sampled vertices coincide. The expected
indicators of a cycle on positions $1,2,3,4$, of the three edges
$12,34,56$, and of the edge $12$ are respectively
$t(C_4,G)$, $\rho^3$, and $\rho$. Averaging over all $6!$ permutations
of the positions gives the target coefficient associated with each host.

Each certificate term is an averaged square. For rootless terms, two
independent triples realize a product of identical flag functions.
For two-root terms, condition on the ordered root and its adjacency;
the two disjoint pairs of extension positions are independent, so the
conditional expectation of the product is a square. For four-root
terms, canonically order the root by its adjacency mask and condition
on it. The two leaf positions are again independent and identically
distributed. Canonicalization depends only on the root, not on either
leaf. Multiplication by the indicator of the specified root type
preserves nonnegativity. This reasoning is valid with replacement and
does not require the original graph to tend to a graphon limit.

All weights in the certificate are positive. Consequently its expected
sum of squares is nonnegative. Coefficientwise domination on every
triangle-free host implies that the target expectation is also
nonnegative, which is exactly \eqref{eq:c4}.

The independent checker can be run as follows:
\begin{lstlisting}
python -S src/verify_four_cycle.py
\end{lstlisting}
It uses only the standard library and does not import the original
verifier, NetworkX, or a table of unlabelled graphs. Passing this check
supports the cited four-cycle inequality; it does not validate unrelated
claims in the source repository.

\section{Limitations and next questions}
The bound in Theorem~\ref{thm:main} exceeds the conjectured value by
$0.00543n^2$. Neither a proof of the conjecture nor a counterexample is
provided. The certificates establish the numerical ingredients of the
written argument; they do not replace an independent review of its
mathematical interpretation. The selection of density ranges and the
parameter $19/20$ may permit further optimization, but merely optimizing
these constants has not been shown to reach $1/50$.

\section*{Computational materials and AI use}
The manuscript, source and computational materials are available at
\url{https://github.com/Komeiji-Shiki/erdos-128-sparse-halves}.
File paths refer to that repository or the unpacked review archive.

AI assistance in Codex was used for literature searches, mathematical
exploration and derivation, programming, drafting, and the subsequent
audit. Earlier methods and Sarid's numerical certificate are attributed
where used. The author reviewed the mathematical arguments, the
applicability of the cited results, and the computational checks.

\begin{thebibliography}{9}
\bibitem{bcl}
J. Balogh, F. C. Clemen and B. Lidick\'y,
\emph{Max Cuts in Triangle-free Graphs}, extended abstract, 2021.
\href{https://arxiv.org/abs/2103.14179v1}{arXiv:2103.14179v1}.

\bibitem{razborov}
A. A. Razborov, \emph{More about sparse halves in triangle-free graphs},
Sbornik: Mathematics \textbf{213} (2022), no.~1, 109--128.
\href{https://doi.org/10.1070/SM9615}{doi:10.1070/SM9615}.
Preprint: \href{https://arxiv.org/abs/2104.09406v2}{arXiv:2104.09406v2}.

\bibitem{sarid}
A. Sarid, \emph{Sparse halves in triangle-free graphs: a bound of
$131/5000$}, preprint dated 4 August 2026, with computational files.
\href{https://github.com/aimir/erdos-128-sparse-halves}{Source repository}.
Accessed 5 September 2026; version pinned in Appendix~\ref{app:c4}.

\bibitem{erdos}
T. Bloom, \emph{Erd\H{o}s Problem \#128}.
\url{https://www.erdosproblems.com/128}.
Accessed 5 September 2026.
\end{thebibliography}
\end{document}
